% Синтетический кейс: класс amsart (теоремы, \maketitle с \address, \subjclass, abstract).
\documentclass[12pt]{amsart}
\usepackage[T2A]{fontenc}
\usepackage[utf8]{inputenc}
\usepackage[russian]{babel}
\usepackage{amssymb}

\theoremstyle{plain}
\newtheorem{theorem}{Теорема}[section]
\newtheorem{lemma}[theorem]{Лемма}
\theoremstyle{definition}
\newtheorem{definition}[theorem]{Определение}
\numberwithin{equation}{section}

\title{Статья в amsart}
\author{Автор Первый}
\address{Университет}
\email{first@example.org}
\subjclass[2020]{94B05}
\keywords{код, вес}

\begin{document}

\begin{abstract}
Краткая аннотация.
\end{abstract}

\maketitle

\section{Введение}

\begin{definition}\label{def:a}
Определение в amsart.
\end{definition}

\begin{theorem}\label{thm:a}
Теорема в amsart с формулой
\begin{equation}\label{eq:a}
  x \in \mathbb{F}_q .
\end{equation}
\end{theorem}

\begin{proof}
По определению~\ref{def:a} и~\eqref{eq:a}.
\end{proof}

\begin{lemma}
Лемма без метки.
\end{lemma}

\end{document}
